% !TeX program = xelatex
\documentclass[12pt,a4paper]{article}
\usepackage[margin=2cm]{geometry}
\usepackage{fontspec}
\usepackage{amsmath,amssymb,enumitem}
\setmainfont{Times New Roman}
\setlength{\parindent}{0pt}
\setlength{\parskip}{6pt}
\begin{document}
\begin{center}
{\Large\bfseries Đề thi Toán vào lớp 10 Hà Nội năm 2018 -- phần Đại số}
\end{center}

\textbf{Câu 1.} Cho hai biểu thức
\[
A=\frac{\sqrt{x}+4}{\sqrt{x}-1},\qquad
B=\frac{3\sqrt{x}+1}{x+2\sqrt{x}-3}-\frac{2}{\sqrt{x}+3},
\quad x\ge 0,\ x\ne 1.
\]
\begin{enumerate}[label=\arabic*)]
  \item Tính giá trị của biểu thức $A$ khi $x=9$.
  \item Chứng minh $B=\dfrac{1}{\sqrt{x}-1}$.
  \item Tìm tất cả giá trị của $x$ để $\dfrac{A}{B}\ge \dfrac{x}{4}+5$.
\end{enumerate}

\textbf{Câu 2.} Một mảnh đất hình chữ nhật có chu vi bằng $28$ mét và độ dài đường chéo bằng $10$ mét. Tính chiều dài và chiều rộng của hình chữ nhật theo đơn vị mét.

\textbf{Câu 3.}
\begin{enumerate}[label=\arabic*)]
  \item Giải hệ phương trình
  \[
  \begin{cases}
  4x-|y+2|=3,\\
  x+2|y+2|=3.
  \end{cases}
  \]
  \item Trong mặt phẳng tọa độ $Oxy$, cho đường thẳng $(d):y=(m+2)x+3$ và parabol $(P):y=x^2$.
  \begin{enumerate}[label=\alph*)]
    \item Chứng minh $(d)$ luôn cắt $(P)$ tại hai điểm phân biệt.
    \item Tìm tất cả giá trị của $m$ để $(d)$ cắt $(P)$ tại hai điểm phân biệt có hoành độ là các số nguyên.
  \end{enumerate}
\end{enumerate}

\textbf{Câu 5.} Tìm giá trị nhỏ nhất của biểu thức
\[
P=\sqrt{1-x}+\sqrt{1+x}+2\sqrt{x}.
\]
\end{document}
